% notes
\providecommand{\redo}[1]{{\color{red} #1}}
\providecommand{\rebuttal}[1]{{#1}}


\providecommand{\nbc}[3]{
 {\colorbox{#3}{\bfseries\sffamily\scriptsize\textcolor{white}{#1}}}
 {\textcolor{#3}{\sf\small\textit{#2}}}
 }
\definecolor{annotcolo}{rgb}{0.1,0.1,0.9}
\providecommand{\nj}[1]{\nbc{NJ}{#1}{annotcolo}}


% config
\providecommand{\smalltextsc}[1]{\textsc{\small #1}}
\providecommand{\scripttextsc}[1]{\textsc{\scriptsize #1}}
\providecommand{\tinytextsc}[1]{\textsc{\tiny #1}}

% sysnames
\providecommand{\livecodebench}{\smalltextsc{LiveCodeBench}}
\providecommand{\lcb}{\smalltextsc{LCB}}
\providecommand{\livecodebenchtitle}{\textsc{LiveCodeBench}}


% general
\providecommand{\llm}{\smalltextsc{LLM}\xspace}
\providecommand{\llms}{\smalltextsc{LLMs}\xspace}
\providecommand{\chainot}{\smalltextsc{COT}\xspace}
\providecommand{\python}{\smalltextsc{Python}\xspace}




%models

%openai
\providecommand{\openai}{\smalltextsc{OpenAI}}
\providecommand{\gpts}{\smalltextsc{GPTs}\xspace}
\providecommand{\codedavinci}{\smalltextsc{Code-Davinci-002}\xspace}
\providecommand{\gptthreefiveturbo}{\smalltextsc{GPT-3.5-turbo}\xspace}
\providecommand{\gptfour}{\smalltextsc{GPT-4}\xspace}
\providecommand{\gptfourturbo}{\smalltextsc{GPT-4-Turbo}\xspace}
\providecommand{\gptfouro}{\smalltextsc{GPT-4-O}\xspace}

%deepmind
\providecommand{\alphacode}{\smalltextsc{AlphaCode}\xspace}
\providecommand{\alphacodeB}[1]{\smalltextsc{AlphaCode-#1B}\xspace}

%anthropic
\providecommand{\claude}{\smalltextsc{Claude}\xspace}
\providecommand{\claudes}{\smalltextsc{Claudes}\xspace}
\providecommand{\claudetwo}{\smalltextsc{Claude-2}\xspace}
\providecommand{\claudethrees}{\smalltextsc{Claude-3s}\xspace}
\providecommand{\claudeinstantone}{\smalltextsc{Claude-Ins-1}\xspace}
\providecommand{\claudeopus}{\smalltextsc{Claude-3-Opus}\xspace}
\providecommand{\claudesonnet}{\smalltextsc{Claude-3-Sonnet}\xspace}

%google
\providecommand{\gemini}{\smalltextsc{Gemini}\xspace}
\providecommand{\geminis}{\smalltextsc{Geminis}\xspace}
\providecommand{\geminipro}{\smalltextsc{Gemini-Pro}\xspace}
\providecommand{\geminiproonefive}{\smalltextsc{Gemini-Pro-1.5}\xspace}
\providecommand{\geminiflash}{\smalltextsc{Gemini-Flash}\xspace}
\providecommand{\geminiultra}{\smalltextsc{Gemini-Ultra}\xspace}

%meta 
\providecommand{\cllamas}{\smalltextsc{CodeLLaMas}\xspace}
\providecommand{\cllama}{\smalltextsc{CodeLLaMa}\xspace}
\providecommand{\cllamains}{\smalltextsc{CodeLLaMa-Ins}\xspace}
\providecommand{\cllamabase}{\smalltextsc{CodeLLaMa-Base}\xspace}
\providecommand{\cllamaB}[1]{\smalltextsc{CL-Ins-#1B}\xspace}
\providecommand{\cllamacommonB}[1]{CL-#1B}
\providecommand{\cllamabaseB}[1]{\smalltextsc{CL-Base-#1B}\xspace}
\providecommand{\cllamaBList}{\smalltextsc{CL-Ins-\{7, 13, 34\}B}\xspace}
\providecommand{\cllamabaseBList}{\smalltextsc{CL-Base-\{7, 13, 34\}B}\xspace}

\providecommand{\llamas}{\smalltextsc{LLaMa-3s}\xspace}

\providecommand{\llamabase}{\smalltextsc{L3-Base}}
\providecommand{\llamabaselist}{\smalltextsc{L3-Base-\{7, 70\}B}\xspace}
\providecommand{\llamainslist}{\smalltextsc{L3-Ins-\{7, 70\}B}\xspace}
\providecommand{\llamaB}[1]{\smalltextsc{L3-Ins-#1B}\xspace}
\providecommand{\llamabaseB}[1]{\smalltextsc{L3-Base-#1B}\xspace}

%finetunes
\providecommand{\wizardlong}{\smalltextsc{WizardCoder}\xspace}
\providecommand{\wizard}{\smalltextsc{WC-34B}\xspace}
\providecommand{\wizardtwo}{\smalltextsc{WC-33B}\xspace}
\providecommand{\phind}{\smalltextsc{Phind-34B}\xspace}
\providecommand{\magicoders}{\smalltextsc{MagiCoders}\xspace}
\providecommand{\magicoder}[1]{\smalltextsc{MC-#1B}\xspace}
\providecommand{\magicoderBLit}{\smalltextsc{MC-\{6.7, 7\}B}\xspace}

%deepseek
\providecommand{\deepseek}{\smalltextsc{DeepSeek}\xspace}
\providecommand{\deepseekbase}{\smalltextsc{DeepSeek-Base}\xspace}
\providecommand{\deepseekinstruct}{\smalltextsc{DeepSeek-Instruct}\xspace}
\providecommand{\deepseeks}{\smalltextsc{DeepSeeks}\xspace}
\providecommand{\deepseekcode}[1]{\smalltextsc{DS}-#1B\xspace}
\providecommand{\deepseekcodeins}{\smalltextsc{DS-Ins}\xspace}
\providecommand{\deepseekcodeB}[1]{\smalltextsc{DS-Ins-#1B}\xspace}
\providecommand{\deepseekbaseB}[1]{\smalltextsc{DS-Base-#1B}\xspace}
\providecommand{\deepseekcodeBList}{\smalltextsc{DS-Ins-\{1.3, 6.7, 33\}B}\xspace}
\providecommand{\deepseekcodebaseBList}{\smalltextsc{DS-Base-\{1.3, 6.7, 33\}B}\xspace}

% codeqwen
\providecommand{\codeqwen}{\smalltextsc{CodeQwen}\xspace}
\providecommand{\codeqwenchat}{\smalltextsc{CodeQwen-Chat}\xspace}

%mistral

\providecommand{\mixtral}{\smalltextsc{Mixtral}\xspace}
\providecommand{\codestral}{\smalltextsc{Codestral}\xspace}
\providecommand{\mistrallarge}{\smalltextsc{Mistral-L}\xspace}
\providecommand{\mistral}{\smalltextsc{Mistral}\xspace}

% bigcode

\providecommand{\starcoder}{\smalltextsc{StarCoder}\xspace}
\providecommand{\starcodertwo}{\smalltextsc{StarCoder2}\xspace}
\providecommand{\starcodertwobase}{\smalltextsc{StarCoder2-Base}\xspace}
\providecommand{\sctwoinsB}[1]{\smalltextsc{SC2-#1B}\xspace}
\providecommand{\sctwobaseB}[1]{\smalltextsc{SC2-Base-#1B}\xspace}
\providecommand{\sctwobaseBList}{\smalltextsc{SC2-Base-\{3,7,15\}B}\xspace}

% cohere
\providecommand{\commandr}{\smalltextsc{CMD-R}\xspace}
\providecommand{\commandrplus}{\smalltextsc{CMD-R+}\xspace}


% datasets
\providecommand{\thestack}{\smalltextsc{Stack}\xspace}
\providecommand{\apps}{\smalltextsc{APPS}\xspace}
\providecommand{\contests}{\smalltextsc{Code-Contests}\xspace}
\providecommand{\humaneval}{\smalltextsc{HumanEval}}
\providecommand{\humanevalplus}{\smalltextsc{HumanEval+}}
\providecommand{\mbpp}{\smalltextsc{MBPP}}
\providecommand{\cruxeval}{\smalltextsc{CRUXEval}}
\providecommand{\codet}{\smalltextsc{CodeT}}
\providecommand{\xcode}{\smalltextsc{xCodeEval}\xspace}
\providecommand{\codescope}{\smalltextsc{CodeScope}\xspace}
\providecommand{\taco}{\smalltextsc{TACO}\xspace}
\providecommand{\usacobench}{\smalltextsc{USCAOBench}\xspace}
%coding websites
\providecommand{\leetcode}{\smalltextsc{LeetCode}\xspace}
\providecommand{\atcoder}{\smalltextsc{AtCoder}\xspace}
\providecommand{\codeforces}{\smalltextsc{CodeForces}\xspace}
\providecommand{\html}{\smalltextsc{HTML}\xspace}

\providecommand{\easy}{\smalltextsc{Easy}\xspace}
\providecommand{\medium}{\smalltextsc{Medium}\xspace}
\providecommand{\hard}{\smalltextsc{Hard}\xspace}

%website scrapes
\providecommand{\problem}{\smalltextsc{$P$}\xspace}
\providecommand{\solution}{\smalltextsc{$S$}\xspace}
\providecommand{\tests}{\smalltextsc{$T$}\xspace}
\providecommand{\testspublic}{\smalltextsc{$T_{p}$}\xspace}
\providecommand{\testshidden}{\smalltextsc{$T_{h}$}\xspace}
\providecommand{\contestdate}{\smalltextsc{$D$}\xspace}

% metrics
\providecommand{\passmetric}[1]{\smalltextsc{Pass@}#1\xspace}



% %%%%
\lstset{
  basicstyle=\ttfamily\small,
  columns=fullflexible,
  % frame=single,
  breaklines=true,
  postbreak=\mbox{$\hookrightarrow$\space},
  rulecolor=\color{black},
}
\definecolor{codegreen}{rgb}{0,0.6,0}
\lstdefinestyle{mystyle}{  
    commentstyle=\color{codegreen},
    keywordstyle=\color{blue},
    basicstyle=\ttfamily\small,
    breakatwhitespace=false,        
    breaklines=true,                 
    captionpos=b,                    
    keepspaces=true,                 
    showspaces=false,                
    showstringspaces=false,
    showtabs=false,                  
    tabsize=2
}
\lstset{style=mystyle}

% if you use cleveref..

% %%%%%%%%%%%%%%%%%%%%%%%%%%%%%%%%
% % THEOREMS
% %%%%%%%%%%%%%%%%%%%%%%%%%%%%%%%%
% \theoremstyle{plain}
% \newtheorem{theorem}{Theorem}[section]
% \newtheorem{proposition}[theorem]{Proposition}
% \newtheorem{lemma}[theorem]{Lemma}
% \newtheorem{corollary}[theorem]{Corollary}
% \theoremstyle{definition}
% \newtheorem{definition}[theorem]{Definition}
% \newtheorem{assumption}[theorem]{Assumption}
% \theoremstyle{remark}
% \newtheorem{remark}[theorem]{Remark}

% % Todonotes is useful during development; simply uncomment the next line
% %    and comment out the line below the next line to turn off comments

% \input{math_commands.tex}
% \newcommand{\todo}[1]{\textcolor{red}{[TODO: #1]}}
\newcommand{\term}[1]{\emph{#1}}
\definecolor{rewriteblue}{RGB}{0,70,170}
\definecolor{mygreen}{RGB}{226,240,217}
\definecolor{myblue}{RGB}{222,235,247}
\definecolor{myorange}{RGB}{251,229,214}
\definecolor{lightblue}{RGB}{232,242,255}
\definecolor{lightorange}{RGB}{255,241,224}
\definecolor{good}{HTML}{1B7837}
\definecolor{bad}{HTML}{B2182B}
\newcommand{\needsrewrite}[1]{\textcolor{rewriteblue}{[SOURCE-DERIVED; REWRITE: #1]}}
\newenvironment{rewriteblock}{%
  \par\begingroup\color{rewriteblue}%
  \noindent
}{%
  \par\endgroup
}
\newenvironment{draftblock}{%
  \par\noindent
}{%
  \par
}
\newcommand{\sourcepaper}[2]{}

\definecolor{lightgreen}{HTML}{C4FFCF}
\definecolor{codegreen}{rgb}{0,0.6,0}

\newcommand{\red}[1]{\textbf{\textcolor{red}{#1}}}
\newcommand{\green}[1]{\textbf{\textcolor{ForestGreen}{#1}}}
\newcommand{\cmark}{\ding{51}}
\newcommand{\xmark}{\ding{55}}
\newcommand{\benchmark}{\textsc{CRUXEval}\xspace}
\newcommand{\benchmarki}{\textsc{CRUXEval-I}\xspace}
\newcommand{\benchmarko}{\textsc{CRUXEval-O}\xspace}
\newcommand{\codellamamed}{\textsc{Code Llama 13B}\xspace}
\newcommand{\codellamalarge}{\textsc{Code Llama 34B}\xspace}
\providecommand{\multirow}[3]{#3}
\providecommand{\hdashline}[1][]{\midrule}
\NewDocumentCommand{\cdashline}{m o}{}
\providecommand{\makecell}[1]{\begin{tabular}{@{}c@{}}#1\end{tabular}}
\providecommand{\smallsec}[1]{\textbf{#1\,\,}}
\providecommand{\llm}{\textsc{LLM}\xspace}
\providecommand{\llms}{\textsc{LLMs}\xspace}
\providecommand{\lean}{\textsc{Lean}\xspace}
\providecommand{\setstretch}[1]{}
\providecommand{\smallsim}{\mathrel{\sim}}
\providecommand{\subfloat}[2][]{%
  \begin{minipage}{0.49\textwidth}%
  \centering
  #2%
  \if\relax\detokenize{#1}\relax\else\subcaption{#1}\fi
  \end{minipage}%
}

\DeclareMathOperator*{\argmax}{arg\,max}
\DeclareMathOperator*{\argmin}{arg\,min}
\lstdefinelanguage{Lean}{
  morekeywords={theorem,lemma,def,by,have,show,calc,let,fun,match,with,if,then,else,import,open,namespace,end,where,structure,inductive,example},
  sensitive=true,
  morecomment=[l]{--},
  morecomment=[s]{/-}{-/},
  morestring=[b]"
}
\lstdefinelanguage{triton}{
  morekeywords={tl,triton,jit,program_id,load,store,arange,zeros,where,constexpr,block_ptr,dtype},
  sensitive=true,
  morecomment=[l]{\#},
  morestring=[b]"
}
\lstdefinelanguage{CUDA}{
  morekeywords={__global__,__device__,__shared__,__constant__,__host__,__syncthreads,threadIdx,blockIdx,blockDim,gridDim},
  sensitive=true,
  morecomment=[l]{//},
  morecomment=[s]{/*}{*/},
  morestring=[b]",
  morestring=[b]'
}
\lstdefinelanguage{Rust}{
  morekeywords={as,break,const,continue,crate,else,enum,extern,false,fn,for,if,impl,in,let,loop,match,mod,move,mut,pub,ref,return,self,Self,static,struct,super,trait,true,type,unsafe,use,where,while,async,await,dyn},
  sensitive=true,
  morecomment=[l]{//},
  morecomment=[s]{/*}{*/},
  morestring=[b]",
  morestring=[b]'
}
\lstdefinelanguage{Python}{
  morekeywords={and,as,assert,break,class,continue,def,del,elif,else,except,False,finally,for,from,global,if,import,in,is,lambda,None,nonlocal,not,or,pass,raise,return,True,try,while,with,yield},
  sensitive=true,
  morecomment=[l]{\#},
  morecomment=[s]{\"\"\"}{\"\"\"},
  morecomment=[s]{\'\'\'}{\'\'\'},
  morestring=[b]',
  morestring=[b]"
}
\lstdefinelanguage{json}{
  morekeywords={true,false,null},
  sensitive=false,
  morestring=[b]",
  morecomment=[l]{//},
}

\lstdefinestyle{mystyle}{
  commentstyle=\color{codegreen},
  keywordstyle=\color{blue},
  basicstyle=\ttfamily\footnotesize,
  breakatwhitespace=false,
  breaklines=true,
  captionpos=b,
  keepspaces=true,
  showspaces=false,
  showstringspaces=false,
  showtabs=false,
  tabsize=2
}
\lstdefinestyle{boxed}{
  frame=single,
  breaklines=true,
  postbreak=\mbox{$\hookrightarrow$\space},
  commentstyle=\color{codegreen},
  keywordstyle=\color{blue},
  basicstyle=\ttfamily\footnotesize,
  breakatwhitespace=false,
  captionpos=t,
  keepspaces=true,
  showspaces=false,
  showstringspaces=false,
  showtabs=false,
  tabsize=2
}
\lstdefinestyle{specblock}{
  frame=single,
  breaklines=true,
  basicstyle=\ttfamily\scriptsize,
  keepspaces=true,
  showstringspaces=false,
  columns=fullflexible
}
\DefineVerbatimEnvironment{codebox}{Verbatim}{fontsize=\small}
\DefineVerbatimEnvironment{examplecodebox}{Verbatim}{fontsize=\small}
\lstdefinestyle{reasoning}{
  breaklines=true,
  commentstyle=\color{codegreen},
  keywordstyle=\color{blue},
  breakautoindent=0pt,
  breakindent=0pt,
  basicstyle=\itshape\footnotesize,
  breakatwhitespace=true,
  captionpos=t,
  keepspaces=true,
  showspaces=false,
  showstringspaces=false,
  showtabs=false
}
\lstdefinestyle{mystylesmaller}{
  commentstyle=\color{codegreen},
  keywordstyle=\color{blue},
  basicstyle=\ttfamily\scriptsize,
  breaklines=true,
  captionpos=b,
  keepspaces=true,
  showspaces=false,
  showstringspaces=false,
  showtabs=false,
  tabsize=2
}
\lstdefinestyle{mystyletiny}{
  commentstyle=\color{codegreen},
  keywordstyle=\color{blue},
  basicstyle=\ttfamily\tiny,
  breaklines=true,
  captionpos=b,
  keepspaces=true,
  showspaces=false,
  showstringspaces=false,
  showtabs=false,
  tabsize=2
}
\lstdefinestyle{leancompact}{
  language=Lean,
  basicstyle=\fontsize{3}{3.5}\ttfamily,
  breaklines=true,
  captionpos=b,
  keepspaces=true,
  showspaces=false,
  showstringspaces=false,
  showtabs=false,
  frame=none
}
\lstset{
  style=mystyle,
  columns=fullflexible,
  frame=single,
  postbreak=\mbox{$\hookrightarrow$\space},
  rulecolor=\color{black},
  moredelim=[is][\color{red}]{!*}{*!},
  literate={₀}{{$_0$}}1
           {₁}{{$_1$}}1
           {₂}{{$_2$}}1
           {₃}{{$_3$}}1
           {₄}{{$_4$}}1
           {₅}{{$_5$}}1
           {₆}{{$_6$}}1
           {₇}{{$_7$}}1
           {₈}{{$_8$}}1
           {₉}{{$_9$}}1
           {¹}{{$^1$}}1
           {²}{{$^2$}}1
           {³}{{$^3$}}1
           {⁻}{{$^-$}}1
           {·}{{$\cdot$}}1
           {×}{{$\times$}}1
           {∀}{{$\forall$}}1
           {∃}{{$\exists$}}1
           {ℕ}{{$\mathbb{N}$}}1
           {ℚ}{{$\mathbb{Q}$}}1
           {ℝ}{{$\mathbb{R}$}}1
           {ℤ}{{$\mathbb{Z}$}}1
           {∧}{{$\wedge$}}1
           {∨}{{$\vee$}}1
           {∈}{{$\in$}}1
           {∉}{{$\notin$}}1
           {∑}{{$\sum$}}1
           {∫}{{$\int$}}1
           {∣}{{$\mid$}}1
           {≠}{{$\ne$}}1
           {≡}{{$\equiv$}}1
           {≤}{{$\le$}}1
           {≥}{{$\ge$}}1
           {←}{{$\leftarrow$}}1
           {→}{{$\to$}}1
           {↑}{{$\uparrow$}}1
           {⊢}{{$\vdash$}}1
           {⟨}{{$\langle$}}1
           {⟩}{{$\rangle$}}1
           {…}{{\ldots}}1
}

\makeatletter
\def\@captype{lstlisting}
\lst@AddToHook{Init}{\def\@captype{lstlisting}}
\makeatother

\theoremstyle{plain}
\newtheorem{theorem}{Theorem}[chapter]
\newtheorem{lemma}[theorem]{Lemma}
\newtheorem{proposition}[theorem]{Proposition}
\newtheorem{corollary}[theorem]{Corollary}

\theoremstyle{definition}
\newtheorem{definition}[theorem]{Definition}
\newtheorem{example}[theorem]{Example}

\theoremstyle{remark}
\newtheorem{remark}[theorem]{Remark}



\providecommand{\smallsim}{\smallsym{\mathrel}{\sim}}

\makeatletter
\providecommand{\smallsym}[2]{#1{\mathpalette\make@small@sym{#2}}}
\providecommand{\make@small@sym}[2]{%
  \vcenter{\hbox{$\m@th\downgrade@style#1#2$}}%
}
\providecommand{\downgrade@style}[1]{%
  \ifx#1\displaystyle\scriptstyle\else
    \ifx#1\textstyle\scriptstyle\else
      \scriptscriptstyle
  \fi\fi
}
\makeatother


\subsection*{Source Note}
This section is derived from \livecodebench{}. Alex is not first author, so this
material is kept in blue and should be rewritten before final submission.

\subsection{Live Benchmarks as Reporting Infrastructure}
\livecodebench{} is included in this chapter for its field-shaping role. Once
coding models began improving quickly, static benchmarks became less useful for
distinguishing frontier systems and less reliable as uncontaminated measures.
\livecodebench{} made time-windowed, contamination-aware coding evaluation easy
to report and compare.

That infrastructure changed model reporting. A model release could be evaluated
not only on a fixed historical benchmark, but also on recent problems that were
less likely to have appeared in training data. The benchmark's multi-scenario
design also encouraged model developers to report code generation together with
repair, execution, and test-output prediction, bringing evaluation closer to the
components of agentic coding workflows.

\subsection{From Observation to Optimization Target}
The field-shaping effect follows the pattern of this chapter. Once
time-windowed performance became visible, contamination-aware coding ability
became a target. Once self-repair and execution were part of the benchmark,
models and agents could be compared on more than direct synthesis. As a result,
\livecodebench{} became useful as a core benchmark for reasoning-oriented and
reinforcement-learning-oriented model releases.

\subsection{Thesis Takeaway}
\livecodebench{} shows that a benchmark can shape not only the score being
reported, but the behavior the field optimizes. By making freshness,
multi-scenario coding, and model comparisons easier to measure, it helped shift
coding evaluation away from static single-task leaderboards.
